\documentclass[a4paper,12pt]{article}
\usepackage[utf8]{inputenc}
\usepackage[italian]{babel}
\usepackage{listings}
\usepackage{xcolor}
\usepackage{float}
\usepackage{placeins}

\lstset{
    language=C,
    backgroundcolor=\color{lightgray!20},
    frame=single,
    breaklines=true,
    numbers=left,
    numberstyle=\tiny,
    stepnumber=1,
    numbersep=5pt,
    captionpos=b,
    basicstyle=\ttfamily\small,
    keywordstyle=\color{blue}\bfseries,
    commentstyle=\color{green!60!black},
    stringstyle=\color{red},
    showstringspaces=false
}

\lstdefinestyle{output}{
    language={},
    backgroundcolor=\color{lightgray!20},
    frame=single,
    breaklines=true,
    numbers=none,
    captionpos=b,
    basicstyle=\ttfamily\small,
    keywordstyle={},
    commentstyle={},
    stringstyle={},
    showstringspaces=false
}

\title{LabReSiD25 - Hands-On 2}
\author{Davide Ferrara - 518629}
\date{}

\begin{document}

\maketitle

\section*{Traccia}
Usando i Linux Socket scrivere un server bloccante (\texttt{server\_tcp\_block.c}),
un server non bloccante (\texttt{server\_tcp\_nonblock.c}) ed un client (\texttt{client\_tcp.c})
che realizzino una comunicazione basata su TCP.
Il client deve inviare la successione dei numeri di Fibonacci, mentre il server
deve effettuare un controllo sulla successione ed inviare una stringa di
acknowledgment (attenzione, diverso da ACK!).

% ============================================================
\section{Client TCP}

\subsection{Obbiettivo}
Il client genera la sequenza di Fibonacci fino a \texttt{F(n)}, la converte in
stringa (es. \texttt{"0 1 1 2 3 5\textbackslash n"}) e la invia al server via TCP.
Attende la risposta (\texttt{ACK} o \texttt{NACK}) e ripete per \texttt{reps}
volte sulla stessa connessione prima di chiudere.
\newline
Accetta due argomenti da riga di comando: \texttt{n} (indice di Fibonacci) e
\texttt{reps} (numero di invii, default 1).

\subsection{client.c}

\subsubsection{Include e costanti}
\texttt{SERVER\_ADDR} è impostato a \texttt{0.0.0.0} per connettersi al
localhost. \texttt{int\_to\_char()} converte un intero in carattere ASCII
sommando \texttt{'0'}, operazione inversa rispetto al parsing del server.

\begin{lstlisting}[caption={client.c - Include e costanti}]
#include "fibo.h"
#include <arpa/inet.h>
#include <netinet/in.h>
#include <stdio.h>
#include <stdlib.h>
#include <string.h>
#include <sys/socket.h>
#include <unistd.h>

#define SERVER_ADDR "0.0.0.0"
#define PORT 8080

char int_to_char(int n) { return '0' + n; }
\end{lstlisting}
\FloatBarrier

\subsubsection{Costruzione del messaggio}
Genera la sequenza con \texttt{fibo\_sequence()} e la converte in stringa
cifra per cifra usando \texttt{int\_to\_char()}. I numeri sono separati da
spazi e la stringa termina con \texttt{\textbackslash n}.

\begin{lstlisting}[caption={client.c - costruzione messaggio}]
int n = atoi(argv[1]);
int reps = (argc == 3) ? atoi(argv[2]) : 1;

// Genera la sequenza di Fibonacci fino a F(n) e la converte in stringa
// Es. n=5 -> "0 1 1 2 3 5\n"
fibo_sequence_t fs;
fibo_init(&fs);
fibo_sequence(&fs, n);
char msg[1024];
int cur = 0;
for (int i = 0; i < fs.idx; i++) {
  if (i > 0) msg[cur++] = ' ';
  msg[cur++] = int_to_char(fs.array[i]);
}
msg[cur++] = '\n';
msg[cur] = '\0';
\end{lstlisting}
\FloatBarrier

\subsubsection{Connessione al server}
\texttt{socket()} crea il socket TCP. \texttt{connect()} stabilisce la
connessione con il server eseguendo il three-way handshake. La connessione
viene aperta una sola volta prima del loop.

\begin{lstlisting}[caption={client.c - socket() e connect()}]
// SOCK_STREAM crea un socket TCP orientato alla connessione
int client_fd = socket(AF_INET, SOCK_STREAM, 0);
if (client_fd == -1) {
  perror("Errore nella creazione del socket");
  exit(1);
}

// connect() stabilisce la connessione TCP con il server (three-way handshake)
if (connect(client_fd, (struct sockaddr *)&server_addr,
            sizeof(server_addr)) == -1) {
  perror("Errore nella connessione");
  exit(1);
}
\end{lstlisting}
\FloatBarrier

\subsubsection{Loop di invio e ricezione}
Il loop invia il messaggio \texttt{reps} volte sulla stessa connessione e
attende la risposta del server ad ogni iterazione. \texttt{close()} chiude
la connessione solo dopo l'ultimo invio.

\begin{lstlisting}[caption={client.c - loop write/read}]
// Invia il messaggio reps volte sulla stessa connessione e attende ACK/NACK
for (int i = 0; i < reps; i++) {
  if (write(client_fd, msg, strlen(msg)) == -1) {
    perror("Errore nel socket");
    exit(1);
  }

  char buf[1024];
  int r = read(client_fd, buf, sizeof(buf));
  buf[r] = '\0';
  printf("[CLIENT.c] %s:%d: %s", SERVER_ADDR, PORT, buf);

  sleep(1);
}

// Chiude la connessione dopo aver inviato tutti i messaggi
close(client_fd);
\end{lstlisting}
\FloatBarrier

\subsection{Output di esempio}

\begin{lstlisting}[style=output,caption={Output client - ./client 5 3}]
[CLIENT.c] Sequenza di fibonacci: 0 1 1 2 3 5
[CLIENT.c] 0.0.0.0:8080: ACK
[CLIENT.c] 0.0.0.0:8080: ACK
[CLIENT.c] 0.0.0.0:8080: ACK
\end{lstlisting}
\FloatBarrier

% ============================================================
\section{Server TCP Bloccante}

\subsection{Obbiettivo}
Il server TCP bloccante riceve dal client una sequenza di Fibonacci sotto forma
di stringa (es. \texttt{"0 1 1 2 3"}), legge il messaggio dalla socket, effettua
il parsing convertendo ogni carattere nel corrispondente intero e ricostruisce
la sequenza. Se la sequenza è corretta risponde con la stringa \texttt{ACK},
altrimenti con \texttt{NACK}.
\newline
Essendo bloccante, il server non può gestire più di un client alla volta:
le chiamate \texttt{accept()} e \texttt{read()} bloccano il processo finché
non arriva una connessione o dei dati.


\subsection{server\_tcp\_block.c}

\subsubsection{Include e costanti}
Includo le librerie necessarie. Rispetto a UDP aggiungo \texttt{netinet/in.h}
per le strutture TCP.

\begin{lstlisting}[caption={server\_tcp\_block.c - Include}]
#include "fibo.h"
#include <arpa/inet.h>
#include <netinet/in.h>
#include <stdio.h>
#include <stdlib.h>
#include <string.h>
#include <sys/socket.h>
#include <unistd.h>

#define PORT 8080
\end{lstlisting}
\FloatBarrier

\subsubsection{parse\_msg}
Parsa il messaggio ricevuto dal client e popola la struttura \texttt{fibo\_sequence\_t}.
I numeri sono separati da spazi. Caratteri \texttt{\textbackslash n} e \texttt{\textbackslash r} vengono ignorati.
\newline
La conversione da carattere a intero sfrutta la tabella ASCII: le cifre
\texttt{'0'...'9'} sono consecutive (48...57), quindi \texttt{'7' - '0' = 55 - 48 = 7}.

\begin{lstlisting}[caption={server\_tcp\_block.c - parse\_msg()}]
void parse_msg(char *msg, fibo_sequence_t *fs) {
  fibo_init(fs);

  while (*msg != '\0') {
    if (*msg == ' ' || *msg == '\n' || *msg == '\r') {
      msg++;
      continue;
    }

    int n = *msg - '0';
    fibo_push(fs, n);

    msg++;
  }
}
\end{lstlisting}
\FloatBarrier

\subsubsection{Creazione del socket}
Creo un socket di tipo \texttt{SOCK\_STREAM} per TCP. La differenza con UDP
(\texttt{SOCK\_DGRAM}) è fondamentale poichè \texttt{SOCK\_STREAM} garantisce
consegna ordinata e affidabile dei dati tramite connessione.
\newline
\texttt{SO\_REUSEADDR} permette di riutilizzare la porta subito dopo la chiusura
del server, evitando l'errore \texttt{Address already in use} causato dallo stato
\texttt{TIME\_WAIT} di TCP.

\begin{lstlisting}[caption={server\_tcp\_block.c - socket() e SO\_REUSEADDR}]
struct sockaddr_in addr;
addr.sin_family = AF_INET;
addr.sin_port = htons(PORT);
addr.sin_addr.s_addr = INADDR_ANY;

// SOCK_STREAM crea un socket TCP orientato alla connessione
int server_fd = socket(AF_INET, SOCK_STREAM, 0);
if (server_fd == -1) {
  perror("Errore nella creazione del socket");
  exit(1);
}

// Evita "Address already in use" al riavvio: TCP tiene il socket
// in TIME_WAIT per ~2 minuti dopo la chiusura
int opt = 1;
setsockopt(server_fd, SOL_SOCKET, SO_REUSEADDR, &opt, sizeof(opt));
\end{lstlisting}
\FloatBarrier

\subsubsection{bind() e listen()}
\texttt{bind()} associa il socket all'indirizzo e porta configurati.
\texttt{listen()} mette il socket in modalità ascolto con un backlog di 5:
il kernel può tenere in coda fino a 5 connessioni in attesa di essere accettate.
\textbf{Nota:} \texttt{listen()} non è bloccante, ritorna subito.

\begin{lstlisting}[caption={server\_tcp\_block.c - bind() e listen()}]
if (bind(server_fd, (struct sockaddr *)&addr, sizeof(addr)) == -1) {
  perror("Errore nel bind del socket");
  exit(1);
}

// Il backlog di 5 indica quante connessioni il kernel puo' accodare
// mentre il server e' occupato a servire un client
if (listen(server_fd, 5) == -1) {
  perror("Error during listening");
}
\end{lstlisting}
\FloatBarrier

\subsubsection{Loop di accettazione: accept()}
Il loop esterno chiama \texttt{accept()}, che \textbf{blocca} il processo finché
un client non si connette. Quando arriva una connessione, \texttt{accept()} ritorna
un nuovo file descriptor (\texttt{client\_fd}) dedicato esclusivamente a quella
connessione. Il \texttt{server\_fd} originale rimane aperto per accettare altri client.
\texttt{accept()} popola anche la struttura \texttt{client\_addr} con IP e porta del client.

\begin{lstlisting}[caption={server\_tcp\_block.c - accept()}]
struct sockaddr_in client_addr;
socklen_t addrlen = sizeof(client_addr);
char buf[1024];

// Loop esterno: il server non termina mai, accetta un client alla volta
while (1) {
  // accept() blocca finche' un client non si connette.
  // Ritorna un nuovo fd dedicato a quella connessione.
  // Finche' siamo qui dentro, nessun altro client viene accettato.
  int client_fd =
      accept(server_fd, (struct sockaddr *)&client_addr, &addrlen);
  if (client_fd == -1) {
    perror("Could not accept connection");
    continue;
  }
  // loop interno ...
}
\end{lstlisting}
\FloatBarrier

\subsubsection{Loop di lettura: read()}
Il loop interno chiama \texttt{read()}, che \textbf{blocca} finché non arrivano dati.
Il valore di ritorno indica:
\begin{itemize}
    \item \texttt{n > 0}: bytes letti correttamente
    \item \texttt{n == 0}: il client ha chiuso la connessione
    \item \texttt{n == -1}: errore
\end{itemize}

\begin{lstlisting}[caption={server\_tcp\_block.c - read()}]
// Loop interno: legge piu' messaggi dalla stessa connessione
// finche' il client non si disconnette
while (1) {
  // read() blocca finche' non arrivano dati.
  // Ritorna 0 quando il client chiude la connessione (EOF).
  ssize_t n = read(client_fd, buf, sizeof(buf));
  if (n == 0) {
    printf("Il Client si e' disconnesso!\n");
    close(client_fd);
    break;
  }
  if (n == -1) {
    printf("Errore del client!\n");
    break;
  }

  char ip_str[INET_ADDRSTRLEN];
  inet_ntop(AF_INET, &client_addr.sin_addr, ip_str, sizeof(ip_str));
  uint16_t port = ntohs(client_addr.sin_port);

  buf[n] = '\0'; // aggiunge terminatore per usare buf come stringa C
  printf("[SERVER_TPC_BLOCK.c] %s:%d: %s", ip_str, port, buf);
  // verifica fibonacci ...
}
\end{lstlisting}
\FloatBarrier

\subsubsection{Verifica Fibonacci e risposta}
Dopo aver ricevuto il messaggio, il server:
\begin{enumerate}
    \item Parsa la sequenza ricevuta dal client con \texttt{parse\_msg()}
    \item Genera la propria sequenza con \texttt{fibo\_sequence()} della stessa lunghezza
    \item Confronta le due sequenze con \texttt{fibo\_compare()}
    \item Invia \texttt{ACK} se corretta, \texttt{NACK} se sbagliata
\end{enumerate}
\begin{lstlisting}[caption={server\_tcp\_block.c - verifica e risposta}]
fibo_sequence_t fs, client_fs;
memset(&fs, 0, sizeof(fibo_sequence_t));
fibo_init(&fs);

// Parsa la stringa ricevuta e genera la sequenza attesa
parse_msg(buf, &client_fs);
fibo_sequence(&fs, client_fs.idx - 1);

if (fibo_compare(&fs, &client_fs) == 0) {
  char *resp = "ACK\n";
  write(client_fd, resp, strlen(resp));
} else {
  char *resp = "NACK\n";
  write(client_fd, resp, strlen(resp));
}
// non chiudiamo: il client puo' inviare altri messaggi
\end{lstlisting}
\FloatBarrier

\subsection{Limitazione: un client alla volta}
Essendo bloccante, il server non può servire più client contemporaneamente.
Mentre il loop interno legge e risponde ai messaggi del client A (porta 57804),
la chiamata \texttt{accept()} non viene mai rieseguita: il client B (porta 57820)
rimane in attesa nel backlog del kernel. Solo dopo che A si disconnette
(\texttt{read()} ritorna 0), il server chiude \texttt{client\_fd}, esce dal loop
interno e torna ad \texttt{accept()}, servendo B.

\subsection{Output di esempio}

\begin{lstlisting}[style=output,caption={Output server - due client sequenziali}]
[TPC SERVER BLOCKING] Server is listening...
[SERVER_TPC_BLOCK.c] 127.0.0.1:57804: 0 1 1 2 3 5
[SERVER_TPC_BLOCK.c] Sent response to 127.0.0.1:57804: ACK
[SERVER_TPC_BLOCK.c] 127.0.0.1:57804: 0 1 1 2 3 5
[SERVER_TPC_BLOCK.c] Sent response to 127.0.0.1:57804: ACK
[SERVER_TPC_BLOCK.c] 127.0.0.1:57804: 0 1 1 2 3 5
[SERVER_TPC_BLOCK.c] Sent response to 127.0.0.1:57804: ACK
Il Client si e' disconnesso!
[SERVER_TPC_BLOCK.c] 127.0.0.1:57820: 0 1 1 2 3 5
[SERVER_TPC_BLOCK.c] Sent response to 127.0.0.1:57820: ACK
[SERVER_TPC_BLOCK.c] 127.0.0.1:57820: 0 1 1 2 3 5
[SERVER_TPC_BLOCK.c] Sent response to 127.0.0.1:57820: ACK
Il Client si e' disconnesso!
\end{lstlisting}
\FloatBarrier

\begin{lstlisting}[style=output,caption={Output client - ./client 5 3}]
[CLIENT.c] Sequenza di fibonacci: 0 1 1 2 3 5
[CLIENT.c] 0.0.0.0:8080: ACK
[CLIENT.c] 0.0.0.0:8080: ACK
[CLIENT.c] 0.0.0.0:8080: ACK
\end{lstlisting}
\FloatBarrier

% ============================================================
\section{Server TCP Non Bloccante}

\subsection{Obbiettivo}
Il server TCP non bloccante riceve dal client una sequenza di Fibonacci sotto
forma di stringa (es. \texttt{"0 1 1 2 3"}), legge il messaggio dalla socket,
effettua il parsing convertendo ogni carattere nel corrispondente intero e
ricostruisce la sequenza. Se la sequenza è corretta risponde con la stringa
\texttt{ACK}, altrimenti con \texttt{NACK}.
\newline
A differenza del server bloccante, le chiamate \texttt{accept()} e \texttt{read()}
non bloccano il processo: ritornano immediatamente con \texttt{EAGAIN} se non ci
sono dati disponibili. Il processo continua a girare nel loop invece di dormire,
realizzando una \textbf{attesa attiva} (busy-waiting).

\subsection{server\_tcp\_nonblock.c}

\subsubsection{Include e costanti}
Rispetto al server bloccante aggiungo \texttt{errno.h} per leggere il codice
di errore dopo le chiamate non bloccanti, e \texttt{fcntl.h} per la funzione
\texttt{fcntl()} che imposta il flag \texttt{O\_NONBLOCK}.

\begin{lstlisting}[caption={server\_tcp\_nonblock.c - Include}]
#include "fibo.h"
#include <arpa/inet.h>
#include <errno.h>
#include <fcntl.h>
#include <netinet/in.h>
#include <stdio.h>
#include <stdlib.h>
#include <string.h>
#include <sys/socket.h>
#include <unistd.h>

#define PORT 8080
\end{lstlisting}
\FloatBarrier

\subsubsection{set\_nonblocking()}
\texttt{fcntl()} legge prima i flag attuali del file descriptor con
\texttt{F\_GETFL}, poi li riscrive aggiungendo \texttt{O\_NONBLOCK} con
\texttt{F\_SETFL}. Da quel momento, qualsiasi chiamata bloccante su quel fd
ritorna immediatamente con \texttt{errno = EAGAIN} se non ci sono dati.

\begin{lstlisting}[caption={server\_tcp\_nonblock.c - set\_nonblocking()}]
/*
 * Imposta il flag O_NONBLOCK sul file descriptor fd tramite fcntl().
 * Dopo questa chiamata, accept() e read() su fd ritornano subito con
 * errno = EAGAIN invece di bloccare il processo in attesa di dati.
 */
static void set_nonblocking(int fd) {
  int flags = fcntl(fd, F_GETFL, 0);
  fcntl(fd, F_SETFL, flags | O_NONBLOCK);
}
\end{lstlisting}
\FloatBarrier

\subsubsection{parse\_msg}
Identica al server bloccante: parsa la stringa ricevuta e popola la struttura
\texttt{fibo\_sequence\_t}. I numeri sono separati da spazi e la conversione
sfrutta la tabella ASCII (\texttt{'7' - '0' = 7}).

\begin{lstlisting}[caption={server\_tcp\_nonblock.c - parse\_msg()}]
void parse_msg(char *msg, fibo_sequence_t *fs) {
  fibo_init(fs);

  while (*msg != '\0') {
    if (*msg == ' ' || *msg == '\n' || *msg == '\r') {
      msg++;
      continue;
    }

    int n = *msg - '0';
    fibo_push(fs, n);

    msg++;
  }
}
\end{lstlisting}
\FloatBarrier

\subsubsection{Creazione del socket e O\_NONBLOCK}
La creazione del socket è identica al server bloccante. La differenza è la
chiamata a \texttt{set\_nonblocking(server\_fd)} prima del \texttt{bind()}: da
quel momento \texttt{accept()} non dormirà mai, ma ritornerà immediatamente.

\begin{lstlisting}[caption={server\_tcp\_nonblock.c - socket() e O\_NONBLOCK}]
// SOCK_STREAM crea un socket TCP orientato alla connessione
int server_fd = socket(AF_INET, SOCK_STREAM, 0);
if (server_fd == -1) {
  perror("Errore nella creazione del socket");
  exit(1);
}

// Evita "Address already in use" al riavvio: TCP tiene il socket
// in TIME_WAIT per ~2 minuti dopo la chiusura
int opt = 1;
setsockopt(server_fd, SOL_SOCKET, SO_REUSEADDR, &opt, sizeof(opt));

// Rende il server fd non bloccante: accept() ritorna EAGAIN invece di
// bloccare il processo quando non ci sono connessioni in arrivo
set_nonblocking(server_fd);
\end{lstlisting}
\FloatBarrier

\subsubsection{Loop esterno: accept() con attesa attiva}
Il loop chiama \texttt{accept()} che ritorna subito. Se \texttt{client\_fd < 0}
e \texttt{errno == EAGAIN} significa che non c'è nessun client in coda: il
loop continua immediatamente senza dormire. Appena arriva una connessione,
anche il \texttt{client\_fd} viene reso non bloccante con
\texttt{set\_nonblocking()}.

\begin{lstlisting}[caption={server\_tcp\_nonblock.c - accept() con attesa attiva}]
// Loop esterno con attesa attiva: il processo non si blocca mai,
// gira continuamente controllando se arriva una connessione
while (1) {
  int client_fd =
      accept(server_fd, (struct sockaddr *)&client_addr, &addrlen);

  if (client_fd < 0) {
    if (errno == EWOULDBLOCK || errno == EAGAIN) {
      // Nessuna connessione in arrivo: continua il loop senza bloccarsi
      continue;
    } else {
      perror("Errore accept");
      exit(1);
    }
  }

  printf("[SERVER_TCP_NONBLOCK.c] Nuovo client connesso\n");

  // Rende anche il client fd non bloccante: read() ritornera' EAGAIN
  // invece di bloccare il processo quando non ci sono dati
  set_nonblocking(client_fd);
  // loop interno ...
}
\end{lstlisting}
\FloatBarrier

\subsubsection{Loop interno: read() con attesa attiva}
\texttt{read()} ritorna immediatamente. Il valore di ritorno indica:
\begin{itemize}
    \item \texttt{n > 0}: bytes letti correttamente
    \item \texttt{n == 0}: il client ha chiuso la connessione (EOF)
    \item \texttt{n < 0} con \texttt{errno == EAGAIN}: nessun dato ora, riprova
    \item \texttt{n < 0} altro errore: errore reale, chiudi la connessione
\end{itemize}

\begin{lstlisting}[caption={server\_tcp\_nonblock.c - read() con attesa attiva}]
// Loop interno con attesa attiva: legge piu' messaggi dalla stessa
// connessione finche' il client non si disconnette
while (1) {
  // read() ritorna subito: EAGAIN se non ci sono dati,
  // 0 se il client ha chiuso la connessione (EOF)
  ssize_t n = read(client_fd, buf, sizeof(buf) - 1);

  if (n < 0) {
    if (errno == EWOULDBLOCK || errno == EAGAIN) {
      // Nessun dato disponibile ora: continua il loop senza bloccarsi
      continue;
    }
    perror("Errore read");
    break;
  }

  if (n == 0) {
    printf("Il Client si e' disconnesso!\n");
    close(client_fd);
    break;
  }

  buf[n] = '\0'; // aggiunge terminatore per usare buf come stringa C
  printf("[SERVER_TCP_NONBLOCK.c] %s:%d: %s", ip_str, port, buf);
  // verifica fibonacci ...
}
\end{lstlisting}
\FloatBarrier

\subsubsection{Verifica Fibonacci e risposta}
Identica al server bloccante: parsa la sequenza ricevuta, genera quella attesa
e confronta. Invia \texttt{ACK} se corretta, \texttt{NACK} altrimenti.

\begin{lstlisting}[caption={server\_tcp\_nonblock.c - verifica e risposta}]
fibo_sequence_t fs, client_fs;
memset(&fs, 0, sizeof(fibo_sequence_t));
fibo_init(&fs);

// Parsa la stringa ricevuta e genera la sequenza attesa
parse_msg(buf, &client_fs);
fibo_sequence(&fs, client_fs.idx - 1);

if (fibo_compare(&fs, &client_fs) == 0) {
  char *resp = "ACK\n";
  write(client_fd, resp, strlen(resp));
  printf("[SERVER_TCP_NONBLOCK.c] Sent response to %s:%d: %s", ip_str,
         port, resp);
} else {
  char *resp = "NACK\n";
  write(client_fd, resp, strlen(resp));
  printf("[SERVER_TCP_NONBLOCK.c] Sent response to %s:%d: %s", ip_str,
         port, resp);
}
// non chiudiamo: il client puo' inviare altri messaggi
\end{lstlisting}
\FloatBarrier

\subsection{Limitazione: ancora un client alla volta}
Nonostante le chiamate siano non bloccanti, il comportamento verso i client è
identico al server bloccante: il loop interno gira su \texttt{read()} finché
il client non si disconnette, e solo allora il loop esterno torna ad
\texttt{accept()}. Il client B deve aspettare che il client A si disconnetta.
\newline
La differenza reale è che il processo non \textbf{dorme} mai: spende il 100\%
della CPU girando nel loop invece di andare in sleep come farebbe il server
bloccante. Questa tecnica, chiamata \textbf{attesa attiva} (busy-waiting), è
più reattiva ma inefficiente in termini di consumo CPU.
\newline
È possibile verificarlo con \texttt{top} mentre entrambi i server sono in
esecuzione senza client connessi:

\begin{lstlisting}[style=output,caption={top -p \$(pgrep -f server\_tcp\_nonblock) - confronto CPU}]
top -p $(pgrep -f server_tcp_block),$(pgrep -f server_tcp_nonblock)

  PID USER  PR  NI    VIRT    RES    SHR S  %CPU  %MEM     TIME+ COMMAND
216495 dff  20   0    2556   1460   1364 R  92.2   0.0  0:37.02 server_tcp_nonb
217176 dff  20   0    2556   1460   1364 S   0.0   0.0  0:00.00 server_tcp_bloc
\end{lstlisting}
\FloatBarrier
Il server non bloccante è nello stato \texttt{R} (Running) al 92\% di CPU,
il server bloccante è nello stato \texttt{S} (Sleeping) allo 0\%: il kernel
lo ha sospeso in attesa di una connessione.
\newline

\subsection{Output di esempio}

\begin{lstlisting}[style=output,caption={Output server - due client sequenziali}]
[SERVER_TCP_NONBLOCK.c] Server is listening...
[SERVER_TCP_NONBLOCK.c] Nuovo client connesso
[SERVER_TCP_NONBLOCK.c] 127.0.0.1:40360: 0 1 1 2 3 5
[SERVER_TCP_NONBLOCK.c] Sent response to 127.0.0.1:40360: ACK
[SERVER_TCP_NONBLOCK.c] 127.0.0.1:40360: 0 1 1 2 3 5
[SERVER_TCP_NONBLOCK.c] Sent response to 127.0.0.1:40360: ACK
[SERVER_TCP_NONBLOCK.c] 127.0.0.1:40360: 0 1 1 2 3 5
[SERVER_TCP_NONBLOCK.c] Sent response to 127.0.0.1:40360: ACK
Il Client si e' disconnesso!
[SERVER_TCP_NONBLOCK.c] Nuovo client connesso
[SERVER_TCP_NONBLOCK.c] 127.0.0.1:40372: 0 1 1 2 3 5
[SERVER_TCP_NONBLOCK.c] Sent response to 127.0.0.1:40372: ACK
Il Client si e' disconnesso!
\end{lstlisting}
\FloatBarrier

\begin{lstlisting}[style=output,caption={Output client A - ./client 5 3}]
[CLIENT.c] Sequenza di fibonacci: 0 1 1 2 3 5
[CLIENT.c] 0.0.0.0:8080: ACK
[CLIENT.c] 0.0.0.0:8080: ACK
[CLIENT.c] 0.0.0.0:8080: ACK
\end{lstlisting}
\FloatBarrier

\begin{lstlisting}[style=output,caption={Output client B - ./client 5 1}]
[CLIENT.c] Sequenza di fibonacci: 0 1 1 2 3 5
[CLIENT.c] 0.0.0.0:8080: ACK
\end{lstlisting}
\FloatBarrier


\end{document}
